% University report template for texforge.
% Text in double braces is a placeholder filled in when the project is created;
% every % comment says what to replace with your own content.
% Tectonic is XeTeX: the source is always UTF-8, so no input encoding
% package is loaded here.
\documentclass[a4paper,12pt]{report}

% Document language, taken from the {{language}} placeholder (default english).
\usepackage[{{language}}]{babel}
\usepackage[T1]{fontenc}
\usepackage[lining]{ebgaramond}
\usepackage{ebgaramond-maths}
\usepackage{microtype}
\usepackage[skip=0.55\baselineskip]{parskip}
\usepackage[a4paper,margin=2.5cm,headheight=14pt,headsep=12pt]{geometry}
\usepackage{xcolor}
\usepackage{amsmath}
\usepackage{booktabs}
\usepackage{float}
\usepackage{graphicx}
\usepackage{fancyhdr}
\usepackage{titlesec}
\usepackage{hyperref}

% One warm accent for navigation, with near-black text for comfortable reading.
\definecolor{reportink}{HTML}{1D1510}
\definecolor{reportaccent}{HTML}{87222C}
\definecolor{reportmark}{HTML}{87222C}
% Fixed labels of the cover, in English; each \addto\captions<lang> block
% below renews them for its language.
\newcommand{\lblAuthor}{Author}
\newcommand{\lblAdvisor}{Advisor}

% Language-dependent defaults for cover placeholders.
% Mechanism: \rptpick{<value macro>}{<default macro>} expands to the default
% when the value macro is empty (\ifx#1\empty#2\else#1\fi).
% Explicit values (from template.toml) always win; only empty values fall back
% to the language default defined in \addto\captions<lang>.
% Using \def (not \newcommand) for value macros: \newcommand creates \long
% macros while \empty is non-\long, and \ifx distinguishes them — so only
% \def-defined value macros compare equal to \empty when empty.
\newcommand{\rptpick}[2]{\ifx#1\empty#2\else#1\fi}

\def\rptsubtitle{{{subtitle}}}
\def\rptadvisor{{{advisor}}}
\def\rptfaculty{{{faculty}}}
\def\rptcity{{{city}}}

% Default language values (English) — overridden by \addto\captions<lang> hooks.
\newcommand{\rptdefsubtitle}{A Structured Academic Report}
\newcommand{\rptdefadvisor}{Advisor Name}
\newcommand{\rptdeffaculty}{Faculty of Engineering}
\newcommand{\rptdefcity}{City, Country}

% Language defaults via babel's \captions<lang> hooks.
\addto\captionsenglish{%
  \renewcommand{\rptdefsubtitle}{A Structured Academic Report}%
  \renewcommand{\rptdefadvisor}{Advisor Name}%
  \renewcommand{\rptdeffaculty}{Faculty of Engineering}%
  \renewcommand{\rptdefcity}{City, Country}%
  \renewcommand{\lblAuthor}{Author}%
  \renewcommand{\lblAdvisor}{Advisor}%
}
\addto\captionsspanish{%
  \renewcommand{\rptdefsubtitle}{Un informe académico estructurado}%
  \renewcommand{\rptdefadvisor}{Nombre del director}%
  \renewcommand{\rptdeffaculty}{Facultad de Ingeniería}%
  \renewcommand{\rptdefcity}{Ciudad, País}%
  \renewcommand{\lblAuthor}{Autor}%
  \renewcommand{\lblAdvisor}{Director}%
  % Spanish-specific names for list of tables, table caption, appendix
  \renewcommand{\listtablename}{Índice de tablas}%
  \renewcommand{\tablename}{Tabla}%
  \renewcommand{\appendixname}{Anexo}%
}
\addto\captionsfrench{%
  \renewcommand{\rptdefsubtitle}{Un rapport universitaire structuré}%
  \renewcommand{\rptdefadvisor}{Nom du directeur}%
  \renewcommand{\rptdeffaculty}{Faculté d'ingénierie}%
  \renewcommand{\rptdefcity}{Ville, Pays}%
  \renewcommand{\lblAuthor}{Auteur}%
  \renewcommand{\lblAdvisor}{Directeur}%
}
\addto\captionsgerman{%
  \renewcommand{\rptdefsubtitle}{Ein strukturierter Studienbericht}%
  \renewcommand{\rptdefadvisor}{Name der Betreuung}%
  \renewcommand{\rptdeffaculty}{Fakultät für Ingenieurwissenschaften}%
  \renewcommand{\rptdefcity}{Stadt, Land}%
  \renewcommand{\lblAuthor}{Autor}%
  \renewcommand{\lblAdvisor}{Betreuung}%
}
\addto\captionsitalian{%
  \renewcommand{\rptdefsubtitle}{Una relazione accademica strutturata}%
  \renewcommand{\rptdefadvisor}{Nome del relatore}%
  \renewcommand{\rptdeffaculty}{Facoltà di Ingegneria}%
  \renewcommand{\rptdefcity}{Città, Paese}%
  \renewcommand{\lblAuthor}{Autore}%
  \renewcommand{\lblAdvisor}{Relatore}%
}
\addto\captionsportuguese{%
  \renewcommand{\rptdefsubtitle}{Um relatório acadêmico estruturado}%
  \renewcommand{\rptdefadvisor}{Nome do orientador}%
  \renewcommand{\rptdeffaculty}{Faculdade de Engenharia}%
  \renewcommand{\rptdefcity}{Cidade, País}%
  \renewcommand{\lblAuthor}{Autor}%
  \renewcommand{\lblAdvisor}{Orientador}%
}
\titleformat{\chapter}[display]
{\normalfont\huge\bfseries\color{reportink}}
{\color{reportaccent}\chaptertitlename\ \thechapter}{14pt}{\Huge}
\titlespacing*{\chapter}{0pt}{-12pt}{18pt}
\titleformat{\section}{\normalfont\Large\bfseries\color{reportink}}
{\textcolor{reportaccent}{\thesection}}{0.75em}{}
\titlespacing*{\section}{0pt}{1.3\baselineskip}{0.45\baselineskip}
\setlength{\parindent}{0pt}
\setlength{\emergencystretch}{2em}
\hypersetup{
  colorlinks=true,
  linkcolor=reportaccent,
  urlcolor=reportaccent,
  citecolor=reportaccent,
  breaklinks=true,
  pdftitle={{{title}}},
  pdfauthor={{{author}}}
}

% Running head: chapter mark on the left, page number on the right.
\pagestyle{fancy}
\fancyhf{}
\fancyhead[L]{\nouppercase{\leftmark}}
\fancyhead[R]{\thepage}
\renewcommand{\headrulewidth}{0.4pt}
% Chapter-opening (and front-matter) pages: page number only.
\fancypagestyle{plain}{%
  \fancyhf{}%
  \fancyhead[R]{\thepage}%
  \renewcommand{\headrulewidth}{0pt}%
}

\begin{document}

  \pagenumbering{roman}

  % ---------------------------------------------------------------- title page
  % Replace every placeholder below; it must stay on exactly one page.
  \begin{titlepage}
    \centering
    \vspace*{1.2cm}
    {\scshape\Large {{institution}} \par}
    \vspace{0.35cm}
    {\large \rptpick{\rptfaculty}{\rptdeffaculty} \par}
    \vspace{0.8cm}
    {\color{reportmark}\rule{4cm}{2.4pt}\par}
    \vspace{1.8cm}
    {\huge\bfseries {{title}} \par}
    \vspace{0.5cm}
    {\Large \rptpick{\rptsubtitle}{\rptdefsubtitle} \par}
    \vspace{2.4cm}
    {\large
      \textbf{\lblAuthor:} {{author}}\\[0.35cm]
      \textbf{\lblAdvisor:} \rptpick{\rptadvisor}{\rptdefadvisor}\par}
    \vfill
    % Replace the city value ({{city}} placeholder) with the city of your institution.
    {\large \rptpick{\rptcity}{\rptdefcity} \par}
    \vspace{0.25cm}
    {\large {{date}} \par}
    \vspace*{1.0cm}
  \end{titlepage}

  % -------------------------------------------------------------- front matter
  \chapter*{\abstractname}
  \addcontentsline{toc}{chapter}{\abstractname}
  % Replace this paragraph with a summary of the problem, the method and the
  % main result (four to six lines are enough).
  This report documents the problem addressed, the method followed and the
  main result obtained. It is organised as a formal university report: the
  background fixes the concepts, the methodology describes the procedure, the
  results present the evidence, and the conclusions state the outcome.

  \tableofcontents
  \listoffigures
  \listoftables

  \clearpage
  \pagenumbering{arabic}

  % ------------------------------------------------------------------ chapters
  \chapter{Introduction}
\label{ch:intro}

% Replace this paragraph with the problem your report addresses and why it
% matters. Cite a source with \cite{...} where appropriate.
Many projects start from a problem that is simple to state but difficult to
solve in a repeatable way. This report studies one such problem: it says
what goes wrong today, why an answer matters, and what evidence would count
as a solution.

% Replace this paragraph with the scope: what is in and what is out.
The scope of this work follows the framework described in~\cite{smith2022}
and covers the design, the implementation and the evaluation of a single
approach. Matters outside that boundary are only mentioned for context.

% Replace with a sentence pointing to the rest of the report.
Chapter~\ref{ch:background} fixes the concepts,
Chapter~\ref{ch:methodology} describes the procedure,
Chapter~\ref{ch:results} presents the evidence, and
Chapter~\ref{ch:conclusions} states the outcome.

\section{Objectives}
\label{sec:objectives}

% Replace with the single general objective of the work.
The general objective of this work is to design and evaluate a repeatable
approach to the problem stated above.

% Replace each item with one specific, measurable objective.
The specific objectives are:
\begin{enumerate}
  \item Collect and organise the input data used by the approach.
  \item Implement the steps that turn the input into a reported result.
  \item Evaluate the result against a small set of representative cases.
\end{enumerate}

  \chapter{Background}
\label{ch:background}

% Replace this paragraph with the key concepts and definitions your work
% builds on. Cite \cite{doe2021} where appropriate.
This chapter fixes the vocabulary used throughout the report. A measurement
is a single recorded value, a sample is a finite collection of measurements,
and a summary is one number computed from a sample. A sample method for
structuring this kind of evidence is described in~\cite{doe2021}.
% APA 7 example (only printed when bibstyle = apa7; plain build unchanged).
\ifx\reportbibstyle\apaseven
In text, the same work reads as \citeA{doe2021}.
\fi

% One labelled equation. Replace it if your work uses a different one, and
% keep \label{eq:mean} so the sentence below still resolves.
The summary used in this report is the arithmetic mean of the sample,
\begin{equation}
  \bar{x} = \frac{1}{n} \sum_{i=1}^{n} x_i
  \label{eq:mean}
\end{equation}
where $x_i$ denotes the $i$-th observation and $n$ the number of
observations.

% Replace this paragraph with a sentence tying the equation to your methodology.
Equation~\eqref{eq:mean} defines the central quantity estimated in
Chapter~\ref{ch:methodology} and reported in Chapter~\ref{ch:results}.

  \chapter{Methodology}
\label{ch:methodology}

% Replace this paragraph with an overview of your approach. Cite
% \cite{jones2023} where appropriate.
The approach follows a fixed sequence of steps so that the same run can be
repeated and compared. Each step takes the output of the previous one and
produces a well-defined result, following the structured workflow
pattern~\cite{jones2023}.

% Mermaid system architecture diagram. Replace the flowchart with your own
% system. Keep style= and the caption (wrap a caption containing a comma in
% braces).
\begin{mermaid}[style=technical, width=0.85\linewidth, pos=H, caption={System architecture}]
  flowchart LR
  A[Input data] --> B[Processing]
  B --> C[Model]
  C --> D[Results]
\end{mermaid}

% Replace this paragraph with a sentence linking the diagram to the
% requirements table. Diagrams carry no label, so refer to them by name
% rather than with \ref.
The stages above turn the input data into the reported results, and each
stage is bound to one functional requirement in the table that follows.

% Booktabs requirements table. Replace the rows with your functional
% requirements; keep the caption and the label.
\begin{table}[H]
  \centering
  \caption{Functional requirements}
  \label{tab:requirements}
  \begin{tabular}{@{}llc@{}}
    \toprule
    ID & Requirement & Priority \\
    \midrule
    R1 & Capture the input data & High \\
    R2 & Validate the input & High \\
    R3 & Produce the report & Medium \\
    \bottomrule
  \end{tabular}
\end{table}

Table~\ref{tab:requirements} lists the functional requirements that drive
the design.

  \chapter{Results}
\label{ch:results}

% Replace this paragraph with a summary of your findings. Cite
% \cite{openlab2024} where appropriate.
The evaluation was run on a small public sample data set~\cite{openlab2024}
and produced the values summarised in this chapter. Every number reported
here comes from the same run, so the procedure can be repeated
independently.

% Graphviz measurement workflow diagram. Replace the graph with your own
% workflow. Keep style= and the caption (wrap a caption containing a comma
% in braces).
\begin{graphviz}[style=editorial, width=0.8\linewidth, pos=H, caption={Measurement workflow}]
  digraph G {
    rankdir=LR
    measure -> validate
    validate -> report
  }
\end{graphviz}

% Replace this paragraph with a sentence linking the diagram to the code
% example. Diagrams carry no label, so refer to them by name rather than
% with \ref.
The workflow above measures each case, validates the measurement and
reports it; the listing below performs the normalisation step of that
workflow.

% Python code block (at most fifteen lines, numbers=true). Replace it with
% your own snippet. A listing is not a figure or a table, so it takes no
% label and appears in neither list.
\begin{code}[lang=python, numbers=true]
def normalize(values):
    """Return the values scaled to the unit interval."""
    if not values:
        return []
    low, high = min(values), max(values)
    if high == low:
        return [0.0 for _ in values]
    return [(v - low) / (high - low) for v in values]
\end{code}

% Replace with a short sentence describing what the code does.
The function rescales the measured values so that the smallest becomes zero
and the largest becomes one, which makes the cases comparable.

  \chapter{Conclusions}
\label{ch:conclusions}

% Replace this paragraph with what your results show.
The results show that the approach meets the objectives stated in
Chapter~\ref{ch:intro}: the input is captured, processed and reported in a
repeatable way, and the normalised values are consistent across the cases
evaluated.

% Replace this paragraph with the limitations and the next step.
The evidence covers a single sample data set, so the conclusions hold for
the cases reported here. The next step is to repeat the evaluation on a
larger set of cases and to automate the run.

% Optional next steps: replace the second paragraph above with an enumerate
% of two or three concrete steps if you prefer a list over a closing
% paragraph.


  % --------------------------------------------------------------- bibliography
  % The unnumbered bibliography chapter is listed in the table of contents;
  % \clearpage + \phantomsection put the entry on the bibliography's own page.
  \bibliographystyle{plain}
  \clearpage
  \phantomsection
  \addcontentsline{toc}{chapter}{\bibname}
  \bibliography{bib/references}

  % ------------------------------------------------------------------ appendix
  \appendix
  % The \appendix command is issued in main.tex, so this file starts with the
% chapter itself; it becomes Appendix A of the report.
\chapter{Supplementary Material}
\label{app:supplementary}

% Replace this paragraph with a note on what the appendix contains.
This appendix collects the symbols used in the report, so that the notation
of Equation~\eqref{eq:mean} can be read without searching the chapters.

% Booktabs notation table. Replace the rows with your own symbols; keep the
% caption and the label so the list of tables stays complete.
\begin{table}[H]
  \centering
  \caption{Notation used in the report}
  \label{tab:notation}
  \begin{tabular}{@{}ll@{}}
    \toprule
    Symbol & Meaning \\
    \midrule
    $x_i$ & Observation $i$ of the sample \\
    $\bar{x}$ & Arithmetic mean of the observations \\
    $n$ & Number of observations \\
    \bottomrule
  \end{tabular}
\end{table}

% Replace this with further supplementary material: a longer listing, an
% extra table or the full data of a single case.


\end{document}
